\begin{table}[!htbp]
    \centering
    \caption{Balance of current sample restrictions: ads specification}
    \label{tab:sample_restrictions_ads}
    \scriptsize
    \begin{adjustbox}{width=\textwidth}
    \begin{tabular}{lcccccc}
    \toprule
     & \multicolumn{2}{c}{Batch 1} & \multicolumn{2}{c}{Batch 2} & \multicolumn{2}{c}{Both batches} \\
    \cmidrule(lr){2-3} \cmidrule(lr){4-5} \cmidrule(lr){6-7}
     & Within current percentile sample & Users with >0 baseline posts & Within current percentile sample & Users with >0 baseline posts & Within current percentile sample & Users with >0 baseline posts \\
    \midrule
    Ads treatment & 0.000 & 0.000 & -0.002 & 0.002 & -0.001 & 0.001 \\
     & (0.001) & (0.002) & (0.002) & (0.003) & (0.001) & (0.001) \\
    Pure control mean & 0.912 & 0.794 & 0.886 & 0.788 & 0.906 & 0.792 \\
    Observations & 141641 & 141641 & 68537 & 68537 & 210178 & 210178 \\
    Influencer-count restriction (< 9) & Yes & Yes & Yes & Yes & Yes & Yes \\
    Stratification block FE & Yes & Yes & Yes & Yes & Yes & Yes \\
    Permutation-based SE & No & No & No & No & No & No \\
    \bottomrule
    \end{tabular}
    \end{adjustbox}
    \floatfoot{\textit{Notes:} Columns report the effect of ads treatment on indicators for satisfying the current ads-analysis sample restrictions. The first outcome is the indicator for remaining inside the current percentile-screen sample, and the second outcome indicates whether a user made more than zero baseline posts. The ads sample keeps the same influencer-count restriction as the main regressions, but these two restriction outcomes are not imposed before estimation. Standard errors are heteroskedasticity-robust.}
\end{table}
